\question[7]
\lstinline|heap-two.c| (es.\ 18) è interattivo con i comandi \lstinline|auth STRING| (\lstinline|malloc| di \lstinline|struct auth { char name[32]; int auth; }|, azzerata), \lstinline|reset| (\lstinline|free(auth)|), \lstinline|service STRING| (\lstinline|strdup|), \lstinline|login| (riesce se \lstinline|auth && auth->auth|). Nessun comando imposta mai \lstinline|auth->auth|.
(a) Perché \lstinline|reset| lascia un puntatore ``stantio''? (b) Qual è la sequenza di comandi dell'exploit e perché \lstinline|service| riusa lo stesso chunk? (c) Perché vanno inviati tutti insieme chiudendo poi lo stdin?
\begin{solution}
(a) \lstinline|free(auth)| libera il chunk ma \emph{non azzera} la variabile globale \lstinline|auth|, che continua a puntare alla memoria liberata: puntatore \emph{stantio} (use-after-free).\\
(b) Sequenza: \lstinline|auth A| $\to$ \lstinline|reset| $\to$ \lstinline|service B...| ($\ge 36$ byte) $\to$ \lstinline|login|. \lstinline|service| fa \lstinline|strdup| di una stringa della stessa taglia del chunk appena liberato, che viene \emph{riusato} (riuso \textbf{LIFO} del fast bin / tcache). La stringa di \lstinline|service| riempie il chunk fino all'\textbf{offset 32} (il campo \lstinline|auth->auth|, dopo \lstinline|name[32]|), scrivendovi un valore non nullo. Ora \lstinline|auth| (ancora sullo stesso chunk) e \lstinline|auth->auth| sono entrambi veri $\to$ login riuscito. Nessun indirizzo: è una vulnerabilità \emph{di contenuto}.\\
(c) Il programma gira in un \lstinline|while(1)| e \emph{non} fa \lstinline|fflush()| tra i prompt: su una pipe (non tty) lo stdout è bufferizzato a blocchi. Vanno inviati tutti i comandi in un colpo e poi chiuso lo stdin (\lstinline|shutdown('send')| / EOF), così \lstinline|fgets| ritorna \lstinline|NULL|, \lstinline|main| esce e l'output bufferizzato viene scaricato. Inviandoli riga per riga con \lstinline|recvline()| ci si bloccherebbe.
\end{solution}
